% --- Archon physics formalization source begin ---
% archon:physics
% archon:covers PhyXMiniProblems/problem_phyx_mini_0867.lean
% archon:source-report reports/phyx_mini/problem_phyx_mini_0867.source.json

\chapter{Physics problem phyx\_mini\_0867}
\label{ch:PhyXMiniProblems_problem_phyx_mini_0867}

\paragraph{Problem source.}
\#\# Physical scenario

The figure shows two uniformly charged spheres.

\#\# Question

How much higher is b's potential compared to a?

\#\# Answer choices

- A: A: 1.2V

- B: B: 140V

- C: C: 2000V

- D: D: 2100V

\#\# Figure caption (auxiliary; use the image as primary evidence)

The image depicts a diagram with two circles representing charged objects. 

- The larger circle on the left has a diameter of 60 cm and a charge of 100 nC, indicated by a "+" sign in its center. Point "a" is marked on the edge of this circle. 

- The smaller circle on the right has a diameter of 10 cm and a charge of 25 nC, also marked with a "+" sign in its center. Point "b" is marked on the edge of this circle. 

- The distance between the two circles (from the point on the left circle to the point on the right circle) is 100 cm.

- The image also includes two labeled distances: the diameter of the larger circle ("60 cm") and the smaller circle ("10 cm"). 

- The distances and charges are annotated next to their respective elements.

\#\# Dataset metadata

Electrostatics; Physical Model Grounding Reasoning, Spatial Relation Reasoning

\paragraph{Recorded answer/context.}
D: D: 2100V

\paragraph{Figure/image path.}
/root/proposal\_for\_physic/science-mango/phyx\_mini\_run/phyx\_data/test\_image/867.png

\paragraph{Formalization target.}
create a compiling Lean file with sorry bodies at `PhyXMiniProblems/problem\_phyx\_mini\_0867.lean`. The Lean declarations must preserve the physical quantities, dimensions or dimensional roles, figure labels, governing-law hypotheses, and final relation expressed by this problem.
Use Mathlib/Physlib names found through LeanExplore where available. If a physics API is missing, introduce faithful local abstractions rather than scalar placeholder aliases.

\begin{theorem}[Physics formalization target]
\label{thm:physics:phyx_mini_0867:target}
\lean{PhyXMiniProblems.ProblemPhyXMini0867.problem_phyx_mini_0867}
\uses{def:physics:phyx-mini-0867:phyxminiproblems-problemphyxmini0867-electricpotentialinvolts, def:physics:phyx-mini-0867:phyxminiproblems-problemphyxmini0867-twouniformlychargedspheressetup, def:physics:phyx-mini-0867:phyxminiproblems-problemphyxmini0867-matchesuniformlychargedspherescenario, def:physics:phyx-mini-0867:phyxminiproblems-problemphyxmini0867-matchesreportedspheremeasurements, def:physics:phyx-mini-0867:phyxminiproblems-problemphyxmini0867-matchessuppliedtwospherefigure, def:physics:phyx-mini-0867:phyxminiproblems-problemphyxmini0867-usesfacingsurfacepointgeometry, def:physics:phyx-mini-0867:phyxminiproblems-problemphyxmini0867-hasphysicaltwosphereparameters, def:physics:phyx-mini-0867:phyxminiproblems-problemphyxmini0867-usestextbookcoulombconstant, def:physics:phyx-mini-0867:phyxminiproblems-problemphyxmini0867-satisfiesuniformspherepotentiallawandsuperposition, def:physics:phyx-mini-0867:phyxminiproblems-problemphyxmini0867-potentialdifferencebminusainvolts, def:physics:phyx-mini-0867:phyxminiproblems-problemphyxmini0867-recordeddatasetanswer, def:physics:phyx-mini-0867:phyxminiproblems-problemphyxmini0867-roundstonearestresolution, def:physics:phyx-mini-0867:phyxminiproblems-problemphyxmini0867-isuniqueclosestdisplayedanswer}
The assigned autoformalize agent should translate this physics problem into Lean declarations in the covered file, with theorem and lemma proof bodies written as `by sorry`.
\end{theorem}
\begin{proof}
This is an autoformalization task, not a proof task. Produce statements that can later be proved by the physics prover without weakening the physical model.
\end{proof}

% --- Archon declaration topology begin ---
\section{Declaration topology}
\begin{definition}[electric Potential Dimension]
  \label{def:physics:phyx-mini-0867:phyxminiproblems-problemphyxmini0867-electricpotentialdimension}
  \lean{PhyXMiniProblems.ProblemPhyXMini0867.electricPotentialDimension}
  Electric potential has dimension M L\textasciicircum{}2 T\textasciicircum{}-2 C\textasciicircum{}-1
\end{definition}
\begin{proof}
  This is the declared model definition of electric Potential Dimension.
\end{proof}
\begin{definition}[coulomb Constant Dimension]
  \label{def:physics:phyx-mini-0867:phyxminiproblems-problemphyxmini0867-coulombconstantdimension}
  \lean{PhyXMiniProblems.ProblemPhyXMini0867.coulombConstantDimension}
  Coulomb's constant has dimension M L\textasciicircum{}3 T\textasciicircum{}-2 C\textasciicircum{}-2
\end{definition}
\begin{proof}
  This is the declared model definition of coulomb Constant Dimension.
\end{proof}
\begin{definition}[Length Quantity]
  \label{def:physics:phyx-mini-0867:phyxminiproblems-problemphyxmini0867-lengthquantity}
  \lean{PhyXMiniProblems.ProblemPhyXMini0867.LengthQuantity}
  A nonnegative, unit-independent physical length
\end{definition}
\begin{proof}
  This is the declared model definition of Length Quantity.
\end{proof}
\begin{definition}[Charge Magnitude Quantity]
  \label{def:physics:phyx-mini-0867:phyxminiproblems-problemphyxmini0867-chargemagnitudequantity}
  \lean{PhyXMiniProblems.ProblemPhyXMini0867.ChargeMagnitudeQuantity}
  A nonnegative total electric charge carried by one sphere
\end{definition}
\begin{proof}
  This is the declared model definition of Charge Magnitude Quantity.
\end{proof}
\begin{definition}[Electric Potential Quantity]
  \label{def:physics:phyx-mini-0867:phyxminiproblems-problemphyxmini0867-electricpotentialquantity}
  \lean{PhyXMiniProblems.ProblemPhyXMini0867.ElectricPotentialQuantity}
  \uses{def:physics:phyx-mini-0867:phyxminiproblems-problemphyxmini0867-electricpotentialdimension}
  A signed electric potential, with its physical dimension retained
\end{definition}
\begin{proof}
  This is the declared model definition of Electric Potential Quantity.
\end{proof}
\begin{definition}[Coulomb Constant Quantity]
  \label{def:physics:phyx-mini-0867:phyxminiproblems-problemphyxmini0867-coulombconstantquantity}
  \lean{PhyXMiniProblems.ProblemPhyXMini0867.CoulombConstantQuantity}
  \uses{def:physics:phyx-mini-0867:phyxminiproblems-problemphyxmini0867-coulombconstantdimension}
  A nonnegative dimensionful value of Coulomb's constant
\end{definition}
\begin{proof}
  This is the declared model definition of Coulomb Constant Quantity.
\end{proof}
\begin{definition}[length Readout]
  \label{def:physics:phyx-mini-0867:phyxminiproblems-problemphyxmini0867-lengthreadout}
  \lean{PhyXMiniProblems.ProblemPhyXMini0867.lengthReadout}
  \uses{def:physics:phyx-mini-0867:phyxminiproblems-problemphyxmini0867-lengthquantity}
  Read a physical length in a selected length unit
\end{definition}
\begin{proof}
  This is the declared model definition of length Readout.
\end{proof}
\begin{definition}[length In Meters]
  \label{def:physics:phyx-mini-0867:phyxminiproblems-problemphyxmini0867-lengthinmeters}
  \lean{PhyXMiniProblems.ProblemPhyXMini0867.lengthInMeters}
  \uses{def:physics:phyx-mini-0867:phyxminiproblems-problemphyxmini0867-lengthquantity, def:physics:phyx-mini-0867:phyxminiproblems-problemphyxmini0867-lengthreadout}
  Read a physical length in metres
\end{definition}
\begin{proof}
  This is the declared model definition of length In Meters.
\end{proof}
\begin{definition}[length In Centimeters]
  \label{def:physics:phyx-mini-0867:phyxminiproblems-problemphyxmini0867-lengthincentimeters}
  \lean{PhyXMiniProblems.ProblemPhyXMini0867.lengthInCentimeters}
  \uses{def:physics:phyx-mini-0867:phyxminiproblems-problemphyxmini0867-lengthquantity, def:physics:phyx-mini-0867:phyxminiproblems-problemphyxmini0867-lengthreadout}
  Read a physical length in centimetres
\end{definition}
\begin{proof}
  This is the declared model definition of length In Centimeters.
\end{proof}
\begin{definition}[charge Magnitude In Coulombs]
  \label{def:physics:phyx-mini-0867:phyxminiproblems-problemphyxmini0867-chargemagnitudeincoulombs}
  \lean{PhyXMiniProblems.ProblemPhyXMini0867.chargeMagnitudeInCoulombs}
  \uses{def:physics:phyx-mini-0867:phyxminiproblems-problemphyxmini0867-chargemagnitudequantity}
  Read a positive sphere charge in coherent-SI coulombs
\end{definition}
\begin{proof}
  This is the declared model definition of charge Magnitude In Coulombs.
\end{proof}
\begin{definition}[charge Magnitude In Nanocoulombs]
  \label{def:physics:phyx-mini-0867:phyxminiproblems-problemphyxmini0867-chargemagnitudeinnanocoulombs}
  \lean{PhyXMiniProblems.ProblemPhyXMini0867.chargeMagnitudeInNanocoulombs}
  \uses{def:physics:phyx-mini-0867:phyxminiproblems-problemphyxmini0867-chargemagnitudequantity, def:physics:phyx-mini-0867:phyxminiproblems-problemphyxmini0867-chargemagnitudeincoulombs}
  Read a positive sphere charge in nanocoulombs
\end{definition}
\begin{proof}
  This is the declared model definition of charge Magnitude In Nanocoulombs.
\end{proof}
\begin{definition}[electric Potential In Volts]
  \label{def:physics:phyx-mini-0867:phyxminiproblems-problemphyxmini0867-electricpotentialinvolts}
  \lean{PhyXMiniProblems.ProblemPhyXMini0867.electricPotentialInVolts}
  \uses{def:physics:phyx-mini-0867:phyxminiproblems-problemphyxmini0867-electricpotentialquantity}
  Read an electric potential in coherent-SI volts
\end{definition}
\begin{proof}
  This is the declared model definition of electric Potential In Volts.
\end{proof}
\begin{definition}[coulomb Constant In Newton Meters Squared Per Coulomb Squared]
  \label{def:physics:phyx-mini-0867:phyxminiproblems-problemphyxmini0867-coulombconstantinnewtonmeterssquaredpercoulombsquared}
  \lean{PhyXMiniProblems.ProblemPhyXMini0867.coulombConstantInNewtonMetersSquaredPerCoulombSquared}
  \uses{def:physics:phyx-mini-0867:phyxminiproblems-problemphyxmini0867-coulombconstantquantity}
  Read Coulomb's constant in newton-metres-squared per coulomb-squared
\end{definition}
\begin{proof}
  This is the declared model definition of coulomb Constant In Newton Meters Squared Per Coulomb Squared.
\end{proof}
\begin{definition}[Sphere Label]
  \label{def:physics:phyx-mini-0867:phyxminiproblems-problemphyxmini0867-spherelabel}
  \lean{PhyXMiniProblems.ProblemPhyXMini0867.SphereLabel}
  The two spheres distinguished by size and horizontal placement
\end{definition}
\begin{proof}
  This is the declared model definition of Sphere Label.
\end{proof}
\begin{definition}[Surface Point Label]
  \label{def:physics:phyx-mini-0867:phyxminiproblems-problemphyxmini0867-surfacepointlabel}
  \lean{PhyXMiniProblems.ProblemPhyXMini0867.SurfacePointLabel}
  The two surface points named in the supplied image
\end{definition}
\begin{proof}
  This is the declared model definition of Surface Point Label.
\end{proof}
\begin{definition}[Horizontal Placement]
  \label{def:physics:phyx-mini-0867:phyxminiproblems-problemphyxmini0867-horizontalplacement}
  \lean{PhyXMiniProblems.ProblemPhyXMini0867.HorizontalPlacement}
  Horizontal placement of a sphere in the drawing
\end{definition}
\begin{proof}
  This is the declared model definition of Horizontal Placement.
\end{proof}
\begin{definition}[Horizontal Surface Side]
  \label{def:physics:phyx-mini-0867:phyxminiproblems-problemphyxmini0867-horizontalsurfaceside}
  \lean{PhyXMiniProblems.ProblemPhyXMini0867.HorizontalSurfaceSide}
  Which horizontal side of a circle carries a labelled surface point
\end{definition}
\begin{proof}
  This is the declared model definition of Horizontal Surface Side.
\end{proof}
\begin{definition}[Separation Arrow Role]
  \label{def:physics:phyx-mini-0867:phyxminiproblems-problemphyxmini0867-separationarrowrole}
  \lean{PhyXMiniProblems.ProblemPhyXMini0867.SeparationArrowRole}
  Meaning of the horizontal 100 cm arrow in the primary raster
\end{definition}
\begin{proof}
  This is the declared model definition of Separation Arrow Role.
\end{proof}
\begin{definition}[Sphere Charge Distribution]
  \label{def:physics:phyx-mini-0867:phyxminiproblems-problemphyxmini0867-spherechargedistribution}
  \lean{PhyXMiniProblems.ProblemPhyXMini0867.SphereChargeDistribution}
  Spherical charge-distribution idealization stated by the problem
\end{definition}
\begin{proof}
  This is the declared model definition of Sphere Charge Distribution.
\end{proof}
\begin{definition}[Electrostatic Model]
  \label{def:physics:phyx-mini-0867:phyxminiproblems-problemphyxmini0867-electrostaticmodel}
  \lean{PhyXMiniProblems.ProblemPhyXMini0867.ElectrostaticModel}
  Electrostatic environmental idealization used for the calculation
\end{definition}
\begin{proof}
  This is the declared model definition of Electrostatic Model.
\end{proof}
\begin{definition}[Two Charged Spheres Figure]
  \label{def:physics:phyx-mini-0867:phyxminiproblems-problemphyxmini0867-twochargedspheresfigure}
  \lean{PhyXMiniProblems.ProblemPhyXMini0867.TwoChargedSpheresFigure}
  \uses{def:physics:phyx-mini-0867:phyxminiproblems-problemphyxmini0867-spherelabel, def:physics:phyx-mini-0867:phyxminiproblems-problemphyxmini0867-surfacepointlabel, def:physics:phyx-mini-0867:phyxminiproblems-problemphyxmini0867-horizontalplacement, def:physics:phyx-mini-0867:phyxminiproblems-problemphyxmini0867-horizontalsurfaceside, def:physics:phyx-mini-0867:phyxminiproblems-problemphyxmini0867-separationarrowrole}
  Literal and relational information visible in image 867. The numerical fields are printed unit-labelled readouts, not replacements for physical lengths or charges
\end{definition}
\begin{proof}
  This is the declared model definition of Two Charged Spheres Figure.
\end{proof}
\begin{definition}[Two Uniformly Charged Spheres Setup]
  \label{def:physics:phyx-mini-0867:phyxminiproblems-problemphyxmini0867-twouniformlychargedspheressetup}
  \lean{PhyXMiniProblems.ProblemPhyXMini0867.TwoUniformlyChargedSpheresSetup}
  \uses{def:physics:phyx-mini-0867:phyxminiproblems-problemphyxmini0867-lengthquantity, def:physics:phyx-mini-0867:phyxminiproblems-problemphyxmini0867-chargemagnitudequantity, def:physics:phyx-mini-0867:phyxminiproblems-problemphyxmini0867-electricpotentialquantity, def:physics:phyx-mini-0867:phyxminiproblems-problemphyxmini0867-coulombconstantquantity, def:physics:phyx-mini-0867:phyxminiproblems-problemphyxmini0867-spherelabel, def:physics:phyx-mini-0867:phyxminiproblems-problemphyxmini0867-surfacepointlabel, def:physics:phyx-mini-0867:phyxminiproblems-problemphyxmini0867-spherechargedistribution, def:physics:phyx-mini-0867:phyxminiproblems-problemphyxmini0867-electrostaticmodel, def:physics:phyx-mini-0867:phyxminiproblems-problemphyxmini0867-twochargedspheresfigure}
  Independent physical data for the two-sphere configuration. The contribution of each sphere and the total potential at each labelled point are stored as physical quantities and constrained only by the governing-law assumptions below
\end{definition}
\begin{proof}
  This is the declared model definition of Two Uniformly Charged Spheres Setup.
\end{proof}
\begin{definition}[Matches Uniformly Charged Sphere Scenario]
  \label{def:physics:phyx-mini-0867:phyxminiproblems-problemphyxmini0867-matchesuniformlychargedspherescenario}
  \lean{PhyXMiniProblems.ProblemPhyXMini0867.MatchesUniformlyChargedSphereScenario}
  \uses{def:physics:phyx-mini-0867:phyxminiproblems-problemphyxmini0867-twouniformlychargedspheressetup}
  Qualitative electrostatic and uniform-charge assumptions from the prose
\end{definition}
\begin{proof}
  This is the declared model definition of Matches Uniformly Charged Sphere Scenario.
\end{proof}
\begin{definition}[Matches Reported Sphere Measurements]
  \label{def:physics:phyx-mini-0867:phyxminiproblems-problemphyxmini0867-matchesreportedspheremeasurements}
  \lean{PhyXMiniProblems.ProblemPhyXMini0867.MatchesReportedSphereMeasurements}
  \uses{def:physics:phyx-mini-0867:phyxminiproblems-problemphyxmini0867-lengthincentimeters, def:physics:phyx-mini-0867:phyxminiproblems-problemphyxmini0867-chargemagnitudeinnanocoulombs, def:physics:phyx-mini-0867:phyxminiproblems-problemphyxmini0867-twouniformlychargedspheressetup}
  The physical measurements stated by the problem. These fields contain no potential or answer-choice information
\end{definition}
\begin{proof}
  This is the declared model definition of Matches Reported Sphere Measurements.
\end{proof}
\begin{definition}[Matches Supplied Two Sphere Figure]
  \label{def:physics:phyx-mini-0867:phyxminiproblems-problemphyxmini0867-matchessuppliedtwospherefigure}
  \lean{PhyXMiniProblems.ProblemPhyXMini0867.MatchesSuppliedTwoSphereFigure}
  \uses{def:physics:phyx-mini-0867:phyxminiproblems-problemphyxmini0867-lengthincentimeters, def:physics:phyx-mini-0867:phyxminiproblems-problemphyxmini0867-chargemagnitudeinnanocoulombs, def:physics:phyx-mini-0867:phyxminiproblems-problemphyxmini0867-twouniformlychargedspheressetup}
  Primary-raster evidence, including the interpretation of the 100 cm arrow as center-to-center. This follows the bitmap itself: the arrow endpoints lie above the two sphere centers
\end{definition}
\begin{proof}
  This is the declared model definition of Matches Supplied Two Sphere Figure.
\end{proof}
\begin{definition}[Uses Facing Surface Point Geometry]
  \label{def:physics:phyx-mini-0867:phyxminiproblems-problemphyxmini0867-usesfacingsurfacepointgeometry}
  \lean{PhyXMiniProblems.ProblemPhyXMini0867.UsesFacingSurfacePointGeometry}
  \uses{def:physics:phyx-mini-0867:phyxminiproblems-problemphyxmini0867-lengthinmeters, def:physics:phyx-mini-0867:phyxminiproblems-problemphyxmini0867-twouniformlychargedspheressetup}
  Geometry derived from the spherical diameters and the two facing surface points. In particular the four center-to-point distances are 0.30 m, 0.70 m, 0.95 m, and 0.05 m after applying the measurement data
\end{definition}
\begin{proof}
  This is the declared model definition of Uses Facing Surface Point Geometry.
\end{proof}
\begin{definition}[Has Physical Two Sphere Parameters]
  \label{def:physics:phyx-mini-0867:phyxminiproblems-problemphyxmini0867-hasphysicaltwosphereparameters}
  \lean{PhyXMiniProblems.ProblemPhyXMini0867.HasPhysicalTwoSphereParameters}
  \uses{def:physics:phyx-mini-0867:phyxminiproblems-problemphyxmini0867-lengthinmeters, def:physics:phyx-mini-0867:phyxminiproblems-problemphyxmini0867-chargemagnitudeincoulombs, def:physics:phyx-mini-0867:phyxminiproblems-problemphyxmini0867-coulombconstantinnewtonmeterssquaredpercoulombsquared, def:physics:phyx-mini-0867:phyxminiproblems-problemphyxmini0867-twouniformlychargedspheressetup}
  Positivity and exterior-point conditions needed by the potential law
\end{definition}
\begin{proof}
  This is the declared model definition of Has Physical Two Sphere Parameters.
\end{proof}
\begin{definition}[Uses Textbook Coulomb Constant]
  \label{def:physics:phyx-mini-0867:phyxminiproblems-problemphyxmini0867-usestextbookcoulombconstant}
  \lean{PhyXMiniProblems.ProblemPhyXMini0867.UsesTextbookCoulombConstant}
  \uses{def:physics:phyx-mini-0867:phyxminiproblems-problemphyxmini0867-coulombconstantinnewtonmeterssquaredpercoulombsquared, def:physics:phyx-mini-0867:phyxminiproblems-problemphyxmini0867-twouniformlychargedspheressetup}
  The textbook numerical calibration of Coulomb's constant used for evaluating the multiple-choice data. It contains no potential difference or answer
\end{definition}
\begin{proof}
  This is the declared model definition of Uses Textbook Coulomb Constant.
\end{proof}
\begin{definition}[Satisfies Uniform Sphere Potential Law And Superposition]
  \label{def:physics:phyx-mini-0867:phyxminiproblems-problemphyxmini0867-satisfiesuniformspherepotentiallawandsuperposition}
  \lean{PhyXMiniProblems.ProblemPhyXMini0867.SatisfiesUniformSpherePotentialLawAndSuperposition}
  \uses{def:physics:phyx-mini-0867:phyxminiproblems-problemphyxmini0867-lengthinmeters, def:physics:phyx-mini-0867:phyxminiproblems-problemphyxmini0867-chargemagnitudeincoulombs, def:physics:phyx-mini-0867:phyxminiproblems-problemphyxmini0867-electricpotentialinvolts, def:physics:phyx-mini-0867:phyxminiproblems-problemphyxmini0867-coulombconstantinnewtonmeterssquaredpercoulombsquared, def:physics:phyx-mini-0867:phyxminiproblems-problemphyxmini0867-spherelabel, def:physics:phyx-mini-0867:phyxminiproblems-problemphyxmini0867-twouniformlychargedspheressetup}
  For a uniformly charged sphere, its potential at a point on or outside the sphere equals k Q / r; electrostatic potentials add linearly. This is the general governing law used here, not the requested numerical conclusion
\end{definition}
\begin{proof}
  This is the declared model definition of Satisfies Uniform Sphere Potential Law And Superposition.
\end{proof}
\begin{definition}[potential Difference BMinus AIn Volts]
  \label{def:physics:phyx-mini-0867:phyxminiproblems-problemphyxmini0867-potentialdifferencebminusainvolts}
  \lean{PhyXMiniProblems.ProblemPhyXMini0867.potentialDifferenceBMinusAInVolts}
  \uses{def:physics:phyx-mini-0867:phyxminiproblems-problemphyxmini0867-electricpotentialinvolts, def:physics:phyx-mini-0867:phyxminiproblems-problemphyxmini0867-twouniformlychargedspheressetup}
  How much higher the potential at b is than the potential at a, in volts
\end{definition}
\begin{proof}
  This is the declared model definition of potential Difference BMinus AIn Volts.
\end{proof}
\begin{definition}[Answer Choice]
  \label{def:physics:phyx-mini-0867:phyxminiproblems-problemphyxmini0867-answerchoice}
  \lean{PhyXMiniProblems.ProblemPhyXMini0867.AnswerChoice}
  The four answer labels printed with the problem
\end{definition}
\begin{proof}
  This is the declared model definition of Answer Choice.
\end{proof}
\begin{definition}[displayed Potential Difference In Volts]
  \label{def:physics:phyx-mini-0867:phyxminiproblems-problemphyxmini0867-displayedpotentialdifferenceinvolts}
  \lean{PhyXMiniProblems.ProblemPhyXMini0867.displayedPotentialDifferenceInVolts}
  \uses{def:physics:phyx-mini-0867:phyxminiproblems-problemphyxmini0867-answerchoice}
  Potential difference printed beside each answer choice, in volts
\end{definition}
\begin{proof}
  This is the declared model definition of displayed Potential Difference In Volts.
\end{proof}
\begin{definition}[recorded Dataset Answer]
  \label{def:physics:phyx-mini-0867:phyxminiproblems-problemphyxmini0867-recordeddatasetanswer}
  \lean{PhyXMiniProblems.ProblemPhyXMini0867.recordedDatasetAnswer}
  \uses{def:physics:phyx-mini-0867:phyxminiproblems-problemphyxmini0867-answerchoice}
  The answer label recorded in the source dataset
\end{definition}
\begin{proof}
  This is the declared model definition of recorded Dataset Answer.
\end{proof}
\begin{definition}[Rounds To Nearest Resolution]
  \label{def:physics:phyx-mini-0867:phyxminiproblems-problemphyxmini0867-roundstonearestresolution}
  \lean{PhyXMiniProblems.ProblemPhyXMini0867.RoundsToNearestResolution}
  A value rounds to a display at the stated positive resolution
\end{definition}
\begin{proof}
  This is the declared model definition of Rounds To Nearest Resolution.
\end{proof}
\begin{definition}[Is Unique Closest Displayed Answer]
  \label{def:physics:phyx-mini-0867:phyxminiproblems-problemphyxmini0867-isuniqueclosestdisplayedanswer}
  \lean{PhyXMiniProblems.ProblemPhyXMini0867.IsUniqueClosestDisplayedAnswer}
  \uses{def:physics:phyx-mini-0867:phyxminiproblems-problemphyxmini0867-twouniformlychargedspheressetup, def:physics:phyx-mini-0867:phyxminiproblems-problemphyxmini0867-potentialdifferencebminusainvolts, def:physics:phyx-mini-0867:phyxminiproblems-problemphyxmini0867-answerchoice, def:physics:phyx-mini-0867:phyxminiproblems-problemphyxmini0867-displayedpotentialdifferenceinvolts}
  A choice is uniquely closest to the calculated potential difference
\end{definition}
\begin{proof}
  This is the declared model definition of Is Unique Closest Displayed Answer.
\end{proof}
\begin{lemma}[potential Difference rounding bounds]
  \label{lem:physics:phyx-mini-0867:phyxminiproblems-problemphyxmini0867-potentialdifference-rounding-bounds}
  \lean{PhyXMiniProblems.ProblemPhyXMini0867.potentialDifference_rounding_bounds}
  \uses{def:physics:phyx-mini-0867:phyxminiproblems-problemphyxmini0867-twouniformlychargedspheressetup, def:physics:phyx-mini-0867:phyxminiproblems-problemphyxmini0867-matchesreportedspheremeasurements, def:physics:phyx-mini-0867:phyxminiproblems-problemphyxmini0867-usesfacingsurfacepointgeometry, def:physics:phyx-mini-0867:phyxminiproblems-problemphyxmini0867-hasphysicaltwosphereparameters, def:physics:phyx-mini-0867:phyxminiproblems-problemphyxmini0867-usestextbookcoulombconstant, def:physics:phyx-mini-0867:phyxminiproblems-problemphyxmini0867-satisfiesuniformspherepotentiallawandsuperposition, def:physics:phyx-mini-0867:phyxminiproblems-problemphyxmini0867-potentialdifferencebminusainvolts}
  The sphere-potential law, measurement conversions, and facing-point geometry place V(b) - V(a) within the interval that rounds to 2100 V at 100 V resolution. This is a derived numerical statement, not a premise
\end{lemma}
\begin{proof}
  The conclusion follows from the governing relations and figure readouts named in the statement of potential Difference rounding bounds.
\end{proof}
% --- Archon declaration topology end ---
% --- Archon physics formalization source end ---
